\section{向量：如何用数字表示一个东西？}

假设我们要让计算机记住两条关联：\textbf{法国 $\rightarrow$ 巴黎}，\textbf{日本 $\rightarrow$ 东京}。人类看得懂文字，但公式需要能计算的对象，所以先给每个“检索线索”一个数字表示。

例如，将法国和日本分别编码为：
\begin{equation}
\vect{k}_{\text{法国}}=\begin{bmatrix}1\\0\end{bmatrix},\qquad
\vect{k}_{\text{日本}}=\begin{bmatrix}0\\1\end{bmatrix}.
\end{equation}

这种有序数字列表叫作\textbf{向量（vector）}。$(1,0)$ 是二维向量；真实模型中的向量可以有几十、几百乃至几千维。这里的数字只是教学示意，并非语言模型真实的语义编码。

\begin{center}
\begin{tikzpicture}[scale=1.02,>=Latex]
\draw[->,gray!70] (-.2,0)--(5.0,0) node[below] {$x$};
\draw[->,gray!70] (0,-.2)--(0,3.0) node[left] {$y$};
\draw[->,line width=1.6pt,blue] (0,0)--(4.0,0) node[above right] {法国 $(1,0)$};
\draw[->,line width=1.6pt,orange] (0,0)--(0,2.45) node[right] {日本 $(0,1)$};
\fill[ink] (0,0) circle(1.2pt);\node[below left] at (0,0) {$0$};
\end{tikzpicture}
\end{center}

向量还有一个重要性质：\textbf{长度}。二维向量 $\vect{a}=(a_1,a_2)$ 的长度就是勾股定理：
\begin{equation}
\|\vect{a}\|=\sqrt{a_1^2+a_2^2}.
\end{equation}
例如 $(0.6,0.8)$ 的长度为 $\sqrt{0.6^2+0.8^2}=1$。长度等于 $1$ 的向量叫\textbf{单位向量}，后面很多计算会方便得多。

\begin{idea}
不要把向量想得太复杂。它就是一个可以进行几何和代数运算的\textbf{数字地址}。我们稍后用它代表记忆的 Key（检索线索），用另一个向量代表 Value（记忆内容）。
\end{idea}
